% Appendix C -- Turbo Initech and other Pascal compilers.  A summary, for
% programmers bringing programs to or from other Pascal compilers, of the
% places where the definition in this book differs from what such
% compilers commonly accept.  Every item restates a rule defined (and
% sourced) in the chapter named; no other compiler's behavior is claimed
% beyond "other compilers may".  Sources: ADR-0007 DEC-02 (subset, the
% dialect pinning), DEC-03, DEC-05; os/tps/tps.pas; probe compilations.
\Gates{test-tps-type-fpc, test-compiler, test-seed-fpc-diff, test-pubs-tps}
\chapter{Turbo Initech and other Pascal compilers}
\ix{compatibility}
\ix{porting programs}

Turbo Initech is a deliberately small Pascal.  A program written for
Turbo Initech is in most cases also a correct program for other Pascal
compilers of the same family; the reverse is often not true.  This
appendix lists the differences, with the chapter that defines each rule.

\section{Lexical differences}
\begin{BorTable}{Tokens}{p{40mm}p{58mm}}
\Hd{Turbo Initech} & \Hd{Notes}\EndHd
\rw{integer}, \rw{boolean}, \rw{char}, \rw{true}, \rw{false}, \rw{forward},
  \rw{write}, \rw{writeln}, \rw{ord}, \rw{chr}, \rw{length} are reserved &
  they cannot be redeclared (Chapter 1)\\
\rw{case} is reserved & there is no \rw{case} statement (Chapter 6)\\
decimal integers only & no hexadecimal, no real numbers, no \code{\#}
  character constants (Chapter 1)\\
\code{\{\$} begins a comment & there are no compiler directives
  (Chapter 14)\\
\end{BorTable}

\section{Types}
\begin{itemize}
\item \rw{integer} is 32 bits; arithmetic wraps around without an error
  (Chapter 3).
\item There are no subranges, enumerated types, sets, pointers, real types,
  variant records, typed files or \code{Text} files (Chapter 3).
\item Record fields are only \rw{integer}, \rw{boolean} or \rw{char}; array
  elements are those or a record type (Chapter 3).
\item An array type has no name and an array is never assigned or passed as
  a whole (Chapters 3 and 7).
\item \rw{string} always has a length byte and a capacity of at most 255
  (Chapter 3).
\end{itemize}

\section{Declarations and routines}
\begin{itemize}
\item Every name, including a routine, must be declared before it is used;
  a routine called before its declaration needs \rw{forward}, even where
  another compiler would find the later declaration (Chapter 7).
\item Routines are not nested, and a routine has no local constants or
  types (Chapter 2).
\item Value parameters are only of the simple types; strings and records are
  passed only by \rw{var}, and a \rw{var} string parameter accepts only a
  string of capacity 255 (Chapter 7).
\item Arrays and files are never parameters (Chapter 7).
\item A function returns \rw{integer}, \rw{boolean} or \rw{char}, and its
  value may not be discarded (Chapter 7).
\item A call is always written with parentheses (Chapter 7).
\end{itemize}

\section{Statements and expressions}
\begin{itemize}
\item \rw{and} and \rw{or} always evaluate both operands (Chapter 5).
\item An expression holds at most one relational operator outside
  parentheses (Chapter 5).
\item There are no \rw{case}, \code{goto} or \code{with} statements
  (Chapter 6).
\item An argument of \std{write}{Write} that begins with a string literal is
  that literal alone; there are no field widths (Chapter 6).
\item Indexes are never checked (Chapter 4).
\end{itemize}

\section{Input and output}
\begin{itemize}
\item \std{write}{Write} and \std{writeln}{Writeln} write to COM1, and a line
  ends with a line feed only (Chapter 12).
\item There is no \code{Read}, \code{Readln} or keyboard input; programs read
  files (Chapter 9).
\item Files are untyped and read and written in bytes;
  \std{reset}{Reset} and \std{rewrite}{Rewrite} take a record size of 1, which
  is also the size when it is left out (Chapter 12).
\item A failed file operation ends the program; there is no way to test for
  the error (Chapter 12).
\item There is no \code{Close}; files are closed by the next
  \std{assign}{Assign}, \std{reset}{Reset} or \std{rewrite}{Rewrite}, and by
  InitechDOS when the program ends (Chapter 12).
\end{itemize}

\section{Writing programs for both}
\ix{portable programs}
A program that is to be compiled both by Turbo Initech and by another
Pascal compiler should use only what this book defines, should give the
record size 1 to \std{reset}{Reset} and \std{rewrite}{Rewrite} explicitly,
should not read a function's name inside the function, should declare
every routine before its first call, and should avoid identifiers that the
other compiler reserves.  The programs of the Turbo Initech manuals are
written so.  Each one is compiled by Turbo Initech and run on InitechDOS;
each, except the two that show a file error and a comment that other
compilers read as a directive, is also compiled and run by another Pascal
compiler; and each that uses no files is also compiled and run by the
genesis compiler.  Every run gives the output printed beside the program.
